\section{Introdu\c{c}\~ao}

O rastreamento do c\^ancer do colo do \'utero depende de programas capazes de
identificar altera\c{c}\~oes precursoras antes da progress\~ao da doen\c{c}a. Em
2020, foram estimados aproximadamente 604 mil novos casos e 342 mil \'obitos no
mundo, com maior carga onde vacina\c{c}\~ao, rastreamento e tratamento s\~ao menos
acess\'iveis \cite{singh2023global}. Na citologia convencional, a automa\c{c}\~ao
tem sido estudada em diferentes escalas: classifica\c{c}\~ao de recortes,
segmenta\c{c}\~ao de estruturas, detec\c{c}\~ao de inst\^ancias e an\'alise de campos
ou l\^aminas \cite{jiang2023systematic,fang2024systematic}.

Essas tarefas pressup\~oem formas distintas de anota\c{c}\~ao. Um r\'otulo de
classe informa o que h\'a no recorte; um ponto informa onde uma inst\^ancia est\'a;
uma caixa acrescenta extens\~ao aproximada; uma m\'ascara descreve fronteira.
Logo, transformar um ponto em caixa n\~ao \'e apenas convers\~ao de formato: \'e uma
decis\~ao sobre qual regi\~ao o detector deve aprender. Em campos com muitas
c\'elulas, o tamanho adotado tamb\'em controla a sobreposi\c{c}\~ao entre alvos e
altera a rela\c{c}\~ao entre localiza\c{c}\~ao e contagem.

A CRIC Cervix disponibiliza 400 imagens convencionais de Papanicolau e 11.534
c\'elulas classificadas, cada uma associada a uma coordenada nuclear
\cite{rezende2021cric}. A base permite estudar detec\c{c}\~ao sem recortar
previamente as c\'elulas, mas n\~ao fornece bounding boxes anat\^omicas. Trabalhos
anteriores usaram a CRIC em detec\c{c}\~ao e classifica\c{c}\~ao
\cite{kalbhor2023cervicell}; ainda assim, a sensibilidade do resultado \`a
geometria criada a partir dos pontos permanece uma quest\~ao experimental
n\~ao investigada. A contagem autom\'atica de c\'elulas por imagem \'e uma tarefa
complementar ao diagn\'ostico --- na contagem celular por redes de regress\~ao,
o erro absoluto m\'edio depende tanto do m\'etodo quanto da densidade dos campos
\cite{xie2018fcrn} --- e, como se ver\'a, o tamanho da pseudo-caixa afeta
localiza\c{c}\~ao e contagem de formas distintas.

Este artigo formula o problema como \emph{engenharia do alvo supervisionado}:
n\~ao se prop\~oe uma nova arquitetura nem se usa a categoria Bethesda como sa\'ida.
Todas as coordenadas s\~ao tratadas como inst\^ancias de uma classe \'unica e o
objetivo \'e determinar em que medida um detector treinado com pseudo-caixas
recupera os pontos anotados e preserva a contagem por campo.

Tr\^es perguntas orientam o desenho experimental. \textbf{RQ1 (alvo):} quanto o
tamanho da pseudo-caixa modifica F1, AP50 e erro de contagem?
\textbf{RQ2 (capacidade):} aumentar a escala de YOLO11n para YOLO11s ou YOLO11m
melhora o compromisso entre localiza\c{c}\~ao e contagem?
\textbf{RQ3 (transfer\^encia):} a configura\c{c}\~ao selecionada na valida\c{c}\~ao
mant\'em desempenho no conjunto de teste retido, e como o erro varia com a
densidade celular?

A principal contribui\c{c}\~ao \'e tornar mensur\'avel uma decis\~ao normalmente
escondida na prepara\c{c}\~ao dos dados: a geometria do alvo criada a partir de
pontos. Todos os resultados s\~ao acompanhados de parti\c{c}\~oes por imagem,
limiar escolhido sem acesso ao teste e intervalos de confian\c{c}a calculados por
bootstrap sobre imagens completas.
